\documentclass[11pt,a4paper]{article}
\usepackage[margin=18mm,headheight=15pt]{geometry}
\usepackage{amsmath,booktabs,array,graphicx,microtype,fancyhdr,xcolor,hyperref}
\definecolor{ink}{HTML}{17354A}
\hypersetup{colorlinks=true,urlcolor=ink,pdftitle={Controlled Latent Skill Composition}}
\pagestyle{fancy}\fancyhf{}\fancyhead[L]{\small\color{ink}TwoRoom and Cube}
\fancyhead[R]{\small Frozen-model composition ablations}\fancyfoot[C]{\small\thepage}
\setlength{\parindent}{0pt}\setlength{\parskip}{4pt}\renewcommand{\arraystretch}{1.08}
\setlength{\emergencystretch}{2em}\DeclareMathSizes{10}{10}{7}{7}
\newcommand{\topic}[1]{\par\smallskip\textbf{#1}\quad}
\begin{document}\fontsize{10}{12}\selectfont
\begin{center}{\LARGE\bfseries\color{ink}Controlled Latent Skill Composition\par}
{\large What to match, whether it matters, and how much to optimize\par}
{\small 26 November 2026 | existing frozen models and benchmark contexts}\end{center}

\topic{Common setup}
Each frozen base-plus-adapter actor generates 16 actions from the current numerical
state, its own skill context, base style and private scalar. Inputs have 10/47
state features and 12/20 context features in TwoRoom/Cube. Composition searches
four existing inputs: $u_K,u_W,a_K,a_W$. Actor $j$ receives base code
$\tilde b_j=(\operatorname{sign}(b_1)a_j,b_2)$, preserving route sign and speed.
Each fixed benchmark case supplies the original base style $b$.
With $d=2$ or 5 action channels, the unchanged objective is
\[
 J_H=\underbrace{\frac1{Hd}\sum_{t=1}^H\|A_K(t)-A_W(t)\|_2^2}_{L_{\rm agree}}
 +\frac{\lambda_a}{2}\sum_{j\in\{K,W\}}(a_j-|b_1|)^2
 +\frac{.01}{2}(u_K^2+u_W^2).
\]
$\lambda_a=1$ in TwoRoom and $.0005$ in Cube. The selected release uses the
private Gaussian point-prior; its optional encoder-cycle penalty is disabled.
All networks, coefficients, bounds and inputs retain their existing definitions.
The grid has 90/360 candidates in TwoRoom/Cube, plus the previous solution on
later replans. Four starts use Adam at learning rate $.12$; code magnitudes lie
in $[.01,2.5]$. TwoRoom retains its negative private-code branches; Cube searches
both signs. Every method executes 8 controls and re-observes. TwoRoom issues
$\operatorname{clip}(A_W+A_K-A_B,-1,1)$; Cube issues $A_W$, with its original
gripper threshold. These different executors are preserved.

\topic{Tasks and evidence}
There are 64 existing contexts per category per environment. TwoRoom requires
goal completion, swept-body keep-out and wall safety, and a waypoint visit before
goal entry. Cube requires held waypoint transport, complete-cube keep-out,
supported carrying, no drop/robot-table contact, and released, settled goal
placement for 10 consecutive controls. Easy cases use the supported geometry;
conflict cases tighten competing requirements. In TwoRoom a nominal waypoint
approach can cross the goal early; in Cube larger keep-out spheres sit closer to
the waypoint. Coupled/OOD moves TwoRoom keep-out into the second room, and widens
Cube starts/goals/arm poses while shifting keep-out along transport. All cases
have independent feasibility witnesses, unavailable to the planner.
\textbf{These banks were used in earlier research; this is exploratory reuse.}
After a Colab reset, all six banks were reconstructed to their exact original
file hashes and all evaluations were rerun. Tables use the complete recovery
run. Its uninstrumented reference was checked against the saved original
aggregate results.

\newpage
\section{RQ1: Does the agreement horizon matter?}
Only $H\in\{4,8,12,16\}$ changes; Adam retains 20 updates. The action MSE is
averaged over $H$, and actual execution stays at 8.
@@HORIZON_TABLE@@
\begin{center}\includegraphics[width=\linewidth]{figures/horizon_success.pdf}\end{center}
@@RQ1_FINDING@@
\textit{Table cells are strict joint-success percentages on the same 64 contexts;
figure bars are 95\% Wilson intervals. Constraint/goal rates, horizon-specific
and common-first8 disagreement, stability and timing are in the accompanying CSVs.}
\topic{The same executed window}
The following comparison measures raw disagreement over the common first eight
commands at the first replan, regardless of the optimized horizon. Values are
MSE multiplied by $10^3$. Per-action curves over all16 predicted commands are
provided separately; later prediction agreement is not a physical task certificate.
@@COMMON8_TABLE@@

\newpage
\section{RQ2: Does lower disagreement predict joint success?}
The primary predictor is \emph{negative raw action disagreement after optimization
at the first replan}, before executing any controls. The outcome is actual strict
joint success under the reference $H=8,N=20$. Total $J_H$ is recorded separately.
There is one score per context, not one independent sample per optimizer update.
AUROCs measure ranking across contexts. RQ3 separately tests refinement within
matched contexts.
@@AUROC_TABLE@@
\begin{center}\includegraphics[width=\linewidth]{figures/agreement_predictiveness.pdf}\end{center}
@@RQ2_FINDING@@
Raw MSE, issued-gripper disagreement and executor deformation are separate
diagnostics. Close actor proposals need not certify their own constraints, and
TwoRoom's issued residual need not equal either proposal. Fixed-$u=-.8$
single-adapter controls check one reference behavior, not the existence of a
successful code anywhere in the policy family. Goal, keep-out, waypoint,
wall/early-entry and Cube grasp/drop/table failures are retained separately.
Later replans and fixed elapsed-control landmarks are descriptive: their cohorts
change, and a violation may already have happened. Their AUROCs are not fresh
prospective failure predictors.

\newpage
\section{RQ3: Does gradient refinement help, and what does it cost?}
Only Adam updates $N\in\{0,5,10,20,40\}$ change, with $H=8$. Grid selection,
learning rate, four starts, regularization and execution remain fixed.
$N=0$ selects the exact best grid candidate. A separate legacy-zero diagnostic
retains the original bounded logit/sigmoid initialization and four-start selection
without updates; this separates initialization effects from gradient refinement.
@@STEPS_TABLE@@
\begin{center}\includegraphics[width=\linewidth]{figures/steps_success.pdf}\end{center}
@@RQ3_FINDING@@
\begin{center}\includegraphics[width=\linewidth]{figures/success_vs_latency.pdf}\end{center}
\textit{Random search retains the original grid and adds bounded uniform latent
samples. Sample counts were fixed without outcomes by matching profiled matrix
arithmetic to Adam20, including the grid, generator forwards and input-gradient
backwards. This is a work proxy, not identical wall time. Three random seeds are
averaged within each context before context-bootstrap uncertainty.}

\newpage
\topic{Runtime and computational accounting}
@@WORK_TABLE@@
@@TIME_TABLE@@
The batch1 timing benchmark uses 12 fixed contexts per environment (four per
category), initial and midpoint observed states, three timing repeats, randomized
method order and warmup. Whole contexts are the bootstrap units. Native evaluation
retains TwoRoom batch64 and Cube batch1; native batch time is not presented as
single-agent latency. Runtime tables include grid work, model calls, backward
samples, refinement intervals, cold calls and total rollout wall time. CUDA event
intervals include host launch gaps and are not sums of kernel busy time.
Native rollout wall times also include physics, recording and occasional CPU
checkpoint processing; matched uncontended batch1 measurements support the
inference-speed comparison.

\topic{Native trajectories and failure interpretation}
@@REPRESENTATIVE_NOTE@@
\begin{center}\includegraphics[width=\linewidth]{figures/representative_overview.pdf}\end{center}

\topic{Controls, reproduction and limits}
@@AUDIT_PARAGRAPH@@
All original weights and sources are preserved; planning uses frozen neural
proposals and zero simulator/task-success queries. TwoRoom uses its unchanged
vectorized closed loop, with independent original-native replay of every path;
Cube runs directly in MuJoCo. There is one frozen model pair per environment.
Adam's grid-based planner is deterministic, so changing its unused RNG seed does
not create independent trials. Success intervals use contexts, and AUROC/paired
comparisons use 2,000 context bootstrap resamples (stratified for pooled results).
Sparse-class intervals cannot resolve unseen class variation; one-class AUROCs
are undefined and valid resample counts are supplied. Replans, repeated seeds and native replays do not increase the
number of independent contexts.

The supplied code/configuration manifest, model/task hashes, aggregate tables,
figures and \texttt{REPRODUCE.txt} identify the exact experiment. Native arrays
and individual episode records remain on the GPU host. Raw task banks must be
restored or reproduced to the recorded hashes before rerunning. No neural model
was retrained and no production configuration or public repository was changed.
\end{document}
